\begin{theorem}[轴对称与点到圆的距离]
    此为“将军饮马”问题的变体。将起点关于“河岸”直线作对称点，最短路径问题转化为对称点与“军营”圆心连线的距离，再减去圆的半径。
\end{theorem}
\begin{example}
    在平面直角坐标系中，军营所在区域的边界为 $x^2+y^2=1$ ，河岸所在直线方程为 $x+y=3$ ，将军从点 $A(0,2)$ 处出发，先到河边饮马，然后再返回军营，如果将军只要到达军营所在区域即回到军营，则这个将军所经过的最短路程为（ \hspace{2em} ）
    \begin{tasks}(4)
        \task $\sqrt{5}$
        \task $\sqrt{5}-1$
        \task $\sqrt{10}$
        \task $\sqrt{10}-1$
    \end{tasks}
\end{example}


\begin{theorem}[轴对称与点到圆的最短距离]
    此问题是“将军饮马”模型与点到圆上动点距离模型的综合应用。将定点 $M$ 关于 $x$ 轴作对称点 $M'$，则 $|PM|+|PQ|$ 的最小值为点 $M'$ 到圆 $C$ 上动点 $Q$ 的最短距离，即 $M'$ 与圆心 $C$ 的距离减去半径。
\end{theorem}
\begin{example}
    已知 $Q$ 为圆 $C:(x-2)^2+(y-1)^2=1$ 上一动点，$M(7,4)$ ，点 $P$ 为 $x$ 轴上一动点，则 $|P M|+|P Q|$ 的最小值为 $\qquad$ ．
\end{example}

\vspace{4em}

\begin{theorem}[轴对称求最短路径]
    求动点到两定点距离之和的最小值问题（将军饮马模型），可将其中一个定点关于动点所在的直线作对称，则最小值等于另一个定点与该对称点之间的距离。
\end{theorem}
\begin{example}
    已知圆 $M_1:(x-2)^2+(y-1)^2=2$ ，圆 $M_2: x^2+y^2-2 x+2 y+1=0$ ，点 $P$ 为 $y$ 轴上的动点，则 $|P M_1|+|P M_2|$ 的最小值为（ \hspace{2em} ）
    \begin{tasks}(4)
        \task $3$
        \task $\sqrt{13}$
        \task $\sqrt{10}$
        \task $\sqrt{5}$
    \end{tasks}
\end{example}


\begin{theorem}[圆与圆的位置关系]
    通过比较两圆的圆心距 $d$ 与两圆半径之和 $R+r$、半径之差 $|R-r|$ 的大小，可以判断两圆的位置关系（相离、外切、相交、内切、内含）。
\end{theorem}
\begin{example}
    设圆 $C_1: x^2+y^2=4$ ，圆 $C_2: x^2+y^2-6 x+8=0$ ，则圆 $C_1, ~ C_2$ 的位置关系是（ \hspace{2em} ）
    \begin{tasks}(4)
        \task 内切
        \task 外切
        \task 相交
        \task 相离
    \end{tasks}
\end{example}

\vspace{4em}

\begin{theorem}[两圆公共弦问题]
    两圆方程作差即可得到它们的公共弦所在直线的方程。然后利用圆心到公共弦的距离、半径和半弦长构成的直角三角形，通过勾股定理求出弦长。
\end{theorem}
\begin{example}
    若圆 $x^2+y^2+2 x-4 y-5=0$ 与 $x^2+y^2+2 x-1=0$ 相交于 $A 、 B$ 两点，则公共弦 $A B$ 的长是（ \hspace{2em} ）。
    \begin{tasks}(4)
        \task $1$
        \task $2$
        \task $3$
        \task $4$
    \end{tasks}
\end{example}

\vspace{4em}

\begin{example}
    圆 $C_1: x^2+y^2+2 x-8 y-8=0$ 与圆 $C_2: x^2+y^2+4 x-4 y-4=0$ 的公共弦长为（ \hspace{2em} ）
    \begin{tasks}(4)
        \task $\frac{\sqrt{55}}{5}$
        \task $\frac{2 \sqrt{55}}{5}$
        \task $\frac{3 \sqrt{55}}{5}$
        \task $\frac{4 \sqrt{55}}{5}$
    \end{tasks}
\end{example}

\vspace{4em}

\begin{theorem}[圆的切线方程]
    若已知点 $P(x_0, y_0)$ 在圆 $(x-a)^2+(y-b)^2=r^2$ 上，则过该点的切线方程为 $(x_0-a)(x-a)+(y_0-b)(y-b)=r^2$。
\end{theorem}
\begin{example}
    已知圆 $C:(x-4)^2+(y+5)^2=2$ ，则圆 $C$ 在点 $P(3,-4)$ 处的切线方程为 $\qquad$ ．
\end{example}

\vspace{4em}

\subsection{填空题}
\begin{theorem}[两圆公切线条数]
    两圆有且仅有三条公切线是两圆外切的充要条件。因此，两圆的圆心距等于两圆半径之和。
\end{theorem}
\begin{example}
    若圆 $C_1: x^2+y^2-2 a y=0(a>0)$ 与圆 $C_2: x^2+y^2-4 x+3=0$ 有且仅有三条公切线，则 $a$ 的值为 $\qquad$ ．
\end{example}

\vspace{4em}